\documentclass[11pt]{article}

% Shared notes settings for Elements of AIML lectures (UPES).
% Keep this file free of \documentclass to allow reuse via \input.

\usepackage[utf8]{inputenc}
\usepackage[T1]{fontenc}
\usepackage[a4paper,margin=1in]{geometry}
\usepackage{hyperref}
\usepackage{xurl}
\usepackage{booktabs}
\usepackage{amsmath}
\usepackage{amssymb}
\usepackage{listings}
\usepackage{xcolor}
\usepackage{enumitem}
\usepackage{parskip}

% Code listings (general-purpose).
\lstset{
  basicstyle=\ttfamily\small,
  keywordstyle=\color{blue},
  commentstyle=\color{gray},
  stringstyle=\color{teal},
  showstringspaces=false,
  columns=fullflexible,
  frame=single,
  framerule=0.2pt,
  breaklines=true
}

\setlist[itemize]{topsep=0.3em,itemsep=0.2em}
\setlist[enumerate]{topsep=0.3em,itemsep=0.2em}

% Simple per-section source line.
\newcommand{\notesources}[1]{\par\noindent{\small\color{gray}\textit{Sources:} \sloppy #1\par}}

% Unit-wise source strings for this subject (used by slides/notes).
% Path is relative to the lecture `latex/` folder (the compilation CWD).
\IfFileExists{../../../_shared/aiml-sources.tex}{%
  % Source strings for Elements of AIML (used by slides + notes).

\newcommand{\AIMLRepoURL}{https://github.com/tali7c/Elements-of-AIML}

\newcommand{\AIMLLabSources}{%
Provost and Fawcett, \textit{Data Science for Business}.;%
McKinney, \textit{Python for Data Analysis}.;%
WEKA Documentation (Waikato Environment for Knowledge Analysis), accessed 2026-02-28.;%
Matplotlib and Seaborn documentation, accessed 2026-02-28.%
}

\newcommand{\AIMLUnitISources}{%
Rich and Knight, \textit{Artificial Intelligence}, McGraw Hill, 2017.;%
Russell and Norvig, \textit{Artificial Intelligence: A Modern Approach} (4e), Ch. 1--2 (Introduction, Intelligent Agents).;%
Poole and Mackworth, \textit{Artificial Intelligence: Foundations of Computational Agents} (2e), selected sections.%
}

\newcommand{\AIMLUnitIISources}{%
Russell and Norvig, \textit{Artificial Intelligence: A Modern Approach} (4e), Ch. 7--9 (Logical Agents, First-Order Logic, Inference).;%
Huth and Ryan, \textit{Logic in Computer Science} (2e), selected sections on propositional and first-order logic.;%
Genesereth and Nilsson, \textit{Logical Foundations of Artificial Intelligence}, selected chapters.%
}

\newcommand{\AIMLUnitIIISources}{%
Mueller and Massaron, \textit{Machine Learning for Dummies}, For Dummies, 2016.;%
Mitchell, \textit{Machine Learning}, selected chapters on learning basics and evaluation.;%
Scikit-learn documentation, accessed 2026-02-28.%
}

\newcommand{\AIMLUnitIVSources}{%
Mueller and Massaron, \textit{Machine Learning for Dummies}, For Dummies, 2016.;%
James, Witten, Hastie, and Tibshirani, \textit{An Introduction to Statistical Learning} (2e), selected chapters.;%
Scikit-learn documentation, accessed 2026-02-28.%
}

\newcommand{\AIMLUnitVSources}{%
NIST, \textit{AI Risk Management Framework (AI RMF 1.0)}, 2023.;%
Barocas, Hardt, and Narayanan, \textit{Fairness and Machine Learning} (online book), selected sections.;%
Knaflic, \textit{Storytelling with Data}, selected chapters on communicating results.%
}
%
}{%
  % Faculty copies live under `private/` so their compile CWD is different.
  \IfFileExists{../../../../public/_shared/aiml-sources.tex}{%
    % Source strings for Elements of AIML (used by slides + notes).

\newcommand{\AIMLRepoURL}{https://github.com/tali7c/Elements-of-AIML}

\newcommand{\AIMLLabSources}{%
Provost and Fawcett, \textit{Data Science for Business}.;%
McKinney, \textit{Python for Data Analysis}.;%
WEKA Documentation (Waikato Environment for Knowledge Analysis), accessed 2026-02-28.;%
Matplotlib and Seaborn documentation, accessed 2026-02-28.%
}

\newcommand{\AIMLUnitISources}{%
Rich and Knight, \textit{Artificial Intelligence}, McGraw Hill, 2017.;%
Russell and Norvig, \textit{Artificial Intelligence: A Modern Approach} (4e), Ch. 1--2 (Introduction, Intelligent Agents).;%
Poole and Mackworth, \textit{Artificial Intelligence: Foundations of Computational Agents} (2e), selected sections.%
}

\newcommand{\AIMLUnitIISources}{%
Russell and Norvig, \textit{Artificial Intelligence: A Modern Approach} (4e), Ch. 7--9 (Logical Agents, First-Order Logic, Inference).;%
Huth and Ryan, \textit{Logic in Computer Science} (2e), selected sections on propositional and first-order logic.;%
Genesereth and Nilsson, \textit{Logical Foundations of Artificial Intelligence}, selected chapters.%
}

\newcommand{\AIMLUnitIIISources}{%
Mueller and Massaron, \textit{Machine Learning for Dummies}, For Dummies, 2016.;%
Mitchell, \textit{Machine Learning}, selected chapters on learning basics and evaluation.;%
Scikit-learn documentation, accessed 2026-02-28.%
}

\newcommand{\AIMLUnitIVSources}{%
Mueller and Massaron, \textit{Machine Learning for Dummies}, For Dummies, 2016.;%
James, Witten, Hastie, and Tibshirani, \textit{An Introduction to Statistical Learning} (2e), selected chapters.;%
Scikit-learn documentation, accessed 2026-02-28.%
}

\newcommand{\AIMLUnitVSources}{%
NIST, \textit{AI Risk Management Framework (AI RMF 1.0)}, 2023.;%
Barocas, Hardt, and Narayanan, \textit{Fairness and Machine Learning} (online book), selected sections.;%
Knaflic, \textit{Storytelling with Data}, selected chapters on communicating results.%
}
%
  }{%
    % Source strings for Elements of AIML (used by slides + notes).

\newcommand{\AIMLRepoURL}{https://github.com/tali7c/Elements-of-AIML}

\newcommand{\AIMLLabSources}{%
Provost and Fawcett, \textit{Data Science for Business}.;%
McKinney, \textit{Python for Data Analysis}.;%
WEKA Documentation (Waikato Environment for Knowledge Analysis), accessed 2026-02-28.;%
Matplotlib and Seaborn documentation, accessed 2026-02-28.%
}

\newcommand{\AIMLUnitISources}{%
Rich and Knight, \textit{Artificial Intelligence}, McGraw Hill, 2017.;%
Russell and Norvig, \textit{Artificial Intelligence: A Modern Approach} (4e), Ch. 1--2 (Introduction, Intelligent Agents).;%
Poole and Mackworth, \textit{Artificial Intelligence: Foundations of Computational Agents} (2e), selected sections.%
}

\newcommand{\AIMLUnitIISources}{%
Russell and Norvig, \textit{Artificial Intelligence: A Modern Approach} (4e), Ch. 7--9 (Logical Agents, First-Order Logic, Inference).;%
Huth and Ryan, \textit{Logic in Computer Science} (2e), selected sections on propositional and first-order logic.;%
Genesereth and Nilsson, \textit{Logical Foundations of Artificial Intelligence}, selected chapters.%
}

\newcommand{\AIMLUnitIIISources}{%
Mueller and Massaron, \textit{Machine Learning for Dummies}, For Dummies, 2016.;%
Mitchell, \textit{Machine Learning}, selected chapters on learning basics and evaluation.;%
Scikit-learn documentation, accessed 2026-02-28.%
}

\newcommand{\AIMLUnitIVSources}{%
Mueller and Massaron, \textit{Machine Learning for Dummies}, For Dummies, 2016.;%
James, Witten, Hastie, and Tibshirani, \textit{An Introduction to Statistical Learning} (2e), selected chapters.;%
Scikit-learn documentation, accessed 2026-02-28.%
}

\newcommand{\AIMLUnitVSources}{%
NIST, \textit{AI Risk Management Framework (AI RMF 1.0)}, 2023.;%
Barocas, Hardt, and Narayanan, \textit{Fairness and Machine Learning} (online book), selected sections.;%
Knaflic, \textit{Storytelling with Data}, selected chapters on communicating results.%
}
%
  }%
}


\title{Elements of AIML Lab\\Experiment 03 (After-submission Reference): WEKA for Clustering}
\author{Tofik Ali}
\date{\today}

\begin{document}
\maketitle
\tableofcontents

\section{How to use this reference}
This document is meant to be shared \textbf{after you submit} Experiment~03.
Use it to cross-check that your work demonstrates:
\begin{itemize}[leftmargin=*]
  \item clustering as an \textbf{unsupervised} task,
  \item at least two clusterers compared on the \textbf{same} dataset,
  \item the effect of \textbf{scaling/normalization} on distance-based clustering,
  \item an interpretation of what each cluster might represent.
\end{itemize}

\section{Quick recap: what clustering is (and is not)}
Clustering groups similar instances \textbf{without labels}.
Clusters are \textbf{not guaranteed} to match human categories; they are patterns discovered from the features.

\subsection{Why scaling matters}
Many clustering methods use distances (e.g., Euclidean distance).
If one feature has a much larger range, it dominates the distance and can change the cluster structure.

\section{Worked example in WEKA (end-to-end)}
\subsection{Dataset choice (example)}
If you need a dataset, WEKA often ships sample datasets under \texttt{data/}.
Numeric datasets work well for k-means.
If your dataset includes a label column, you may keep it for \emph{interpretation only} (clustering itself does not use it).

\subsection{Step 1: load data and remove obvious IDs}
\begin{enumerate}[leftmargin=*]
  \item WEKA $\rightarrow$ \textbf{Explorer} $\rightarrow$ \textbf{Preprocess} $\rightarrow$ \textbf{Open file...}
  \item Inspect attributes and remove ID-like columns (if any).
\end{enumerate}

\subsection{Step 2: run k-means once without scaling}
\begin{enumerate}[leftmargin=*]
  \item Go to \textbf{Cluster} tab.
  \item \textbf{Choose} $\rightarrow$ \texttt{SimpleKMeans}.
  \item Set \textbf{NumClusters} (choose a reasonable $k$, e.g., 3 for a first run).
  \item Click \textbf{Start}.
\end{enumerate}
Record:
\begin{itemize}[leftmargin=*]
  \item chosen $k$,
  \item cluster sizes,
  \item centroid summaries (if shown),
  \item any available within-cluster error measure reported by WEKA.
\end{itemize}

\subsection{Step 3: apply scaling and rerun k-means}
\begin{enumerate}[leftmargin=*]
  \item \textbf{Preprocess} tab $\rightarrow$ \textbf{Filter}.
  \item Try \texttt{unsupervised/attribute/Standardize} (or \texttt{Normalize}).
  \item Apply the filter.
  \item Return to \textbf{Cluster} tab and run \texttt{SimpleKMeans} again (same $k$ first).
\end{enumerate}
\textbf{Cross-check:} cluster sizes/centroids often change after scaling if feature ranges differ.

\subsection{Step 4: run a second clusterer (comparison)}
Choose one more:
\begin{itemize}[leftmargin=*]
  \item \texttt{EM} (probabilistic),
  \item \texttt{HierarchicalClusterer} (hierarchical).
\end{itemize}
Run it on the same (preferably scaled) dataset and compare:
\begin{itemize}[leftmargin=*]
  \item cluster sizes,
  \item stability/interpretability,
  \item whether the grouping aligns with your domain intuition.
\end{itemize}

\subsection{Step 5: visualize cluster assignments}
If available, use:
\begin{itemize}[leftmargin=*]
  \item \textbf{Visualize cluster assignments} to view 2D projections.
\end{itemize}
Take at least one screenshot and write 2--3 lines describing what the visualization suggests.

\section{How to write interpretations (examples)}
Good interpretation connects clusters to real meaning:
\begin{itemize}[leftmargin=*]
  \item ``Cluster~1 has high income and low spending score; likely \emph{savers}. Cluster~2 has moderate income and high spending score; likely \emph{active shoppers}.'' (example pattern)
  \item ``After scaling, clusters became more balanced; before scaling, the largest-range feature dominated grouping.''
\end{itemize}

\section{Troubleshooting (common issues and fixes)}
\begin{itemize}[leftmargin=*]
  \item \textbf{Meaningless clusters:} remove IDs, scale numeric features, and re-check $k$.
  \item \textbf{Choosing $k$ randomly:} justify using a simple heuristic (try multiple $k$ values and compare within-cluster error + interpretability).
  \item \textbf{Treating clusters as labels:} clusters are exploratory; report them as patterns, not truth.
\end{itemize}

\section{Report structure (1--2 pages)}
\begin{enumerate}[leftmargin=*]
  \item Dataset (features used, rows/columns, any preprocessing).
  \item Clusterer~1 results (k-means): $k$, cluster sizes, centroid summary highlights.
  \item Scaling comparison: what changed after \texttt{Standardize/Normalize}.
  \item Clusterer~2 results (EM/hierarchical) and comparison.
  \item Interpretation: what each cluster could represent.
\end{enumerate}

\notesources{\AIMLLabSources}
\end{document}

